\index{BPatch_function : : : : : : :!}{

An object of this class represents a function in the application. A BPatch_image object (see description below) can

be used to retrieve a BPatch_function object representing a given function.}



{

std::string getName();}



{

std::string getDemangledName();}



{

std::string getMangledName();}



{

std::string getTypedName();}



{

void getNames(std::vector<std::string> &names);}



{

void getDemangledNames(std::vector<std::string> &names);}



{

void getMangledNames(std::vector<std::string> &names);}



{

void getTypedNames(std::vector<std::string> &names);}



{

Return name(s) of the function. The getName functions return the primary name; this is typically the first symbol we

encounter while parsing the program; getName is an alias for getDemangledName. The getNames functions return all known

names for the function, including any names specified by weak symbols.



{

bool getAddressRange(Dyninst::Address &start,}



{

Dyninst::Address &end)}



{

Returns the bounds of the function; for non-contiguous functions, this is the lowest and highest address of code that

the function includes.}





\bigskip



{

std::vector<BPatch_localVar *> *getParams\index{getParams: : : : : : :!}()}



{

Return a vector of BPatch_localVar snippets that refer to the parameters of this function. The position in the vector

corresponds to the position in the parameter list (starting from zero). The returned local variables can be used to

check the types of functions, and can be used in snippet expressions.}



{

BPatch_type *getReturnType\index{getReturnType: : : : : : :!}()}



{

Return the type of the return value for this function.}



{

BPatch_variableExpr *getFunctionRef()}



{

For platforms with complex function pointers (e.g., 64-bit PPC) this constructs and returns the appropriate descriptor.}



{

std::vector<BPatch_localVar *> *getVars()}



{

Returns a vector of BPatch_localVar objects that contain the local variables in this function. These

BPatch_localVars can be used as parts of snippets in instrumentation. This function requires debug information to be

present in the mutatee. If Dyninst was unable to find any local variables, this function will return an empty vector.

It is up to the user to free the vector returned by this function.}



{

bool isInstrumentable\index{isInstrumentable: : : : : : :!}()}



{

Return true if the function can be instrumented, and false if it cannot. Various conditions can cause a function to be

uninstrumentable. For example, there exists a platform-specific minimum function size beyond which a function cannot

be instrumented.}



{

bool isSharedLib\index{isSharedLib: : : : : : :!}()}



{

This function returns true if the function is defined in a shared library.}



{

BPatch_module *getModule\index{getModule: : : : : : :!}()}



{

Return the module that contains this function. Depending on whether the program was compiled for debugging or the

symbol table stripped, this information may not be available. This function returns NULL if module information was

not found.}



{

char *getModuleName\index{getModuleName: : : : : : :!}(char *name, int maxLen)}



{

Copies the name of the module that contains this function into the buffer pointed to by name. Copies at most maxLen

characters and returns a pointer to name.}



{

enum BPatch_procedureLocation {



{

BPatch_entry,



{

BPatch_exit,



{

BPatch_subroutine,



{

BPatch_locInstruction,}



{

BPatch_locBasicBlockEntry,}



{

BPatch_locLoopEntry,}



{

BPatch_locLoopExit,}



{

BPatch_locLoopStartIter,}



{

BPatch_locLoopStartExit,}



{

BPatch_allLocations }}



{

const std::vector<BPatch_point *> *findPoint\index{findPoint: : : : : : :!}(const

BPatch_procedureLocation loc)}



{

Return the BPatch_point or list of BPatch_points associated with the procedure. It is used to select which type of

points associated with the procedure will be returned. BPatch_entry and BPatch_exit request respectively the entry

and exit points of the subroutine. BPatch_subroutine returns the list of points where the procedure calls other

procedures. If the lookup fails to locate any points of the requested type, NULL is returned.}



{

enum BPatch_opCode { BPatch_opLoad, BPatch_opStore, BPatch_opPrefetch }}



{

std::vector<BPatch_point *> *findPoint(const std::set<BPatch_opCode>&

ops)}



{

std::vector<BPatch_point *> *findPoint(const

BPatch_Set<BPatch_opCode>& ops)\index{findPoint: : : : : : :!}\index{BPatch_opCode: : : : :

: :!}}



{

Return the vector of BPatch_points corresponding to the set of machine instruction types described by the argument.

This version is used primarily for memory access instrumentation. The BPatch_opCode is an enumeration of instruction

types that may be requested: BPatch_opLoad, BPatch_opStore, and BPatch_opPrefetch. Any combination of these may be

requested by passing an appropriate argument set containing the desired types. The instrumentation points created by

this function have additional memory access information attached to them. This allows such points to be used for memory

access specific snippets (e.g. effective address). The memory access information attached is described under Memory

Access classes in section \ref{bkm:Ref160277443}.



{

BPatch_localVar *findLocalVar(const char *name)}



{

Search the function's local variable collection for name. This returns a pointer to the local

variable if a match is found. This function returns NULL if it fails to find any variables.}



{

std::vector<BPatch_variableExpr *> *findVariable(const char * name)}



{

bool findVariable(const char *name, std::vector<BPatch_variableExpr> &vars)}



{

Return a set of variables matching name at the scope of this function. If no variables match in the local scope, then

the global scope will be searched for matches. This function returns NULL if it fails to find any variables.}



{

BPatch_localVar *findLocalParam(const char *name)}



{

Search the function's parameters for a given name. A BPatch_localVar * pointer is returned if a

match is found, and NULL is returned otherwise.}



{

void *getBaseAddr\index{getBaseAddr: : : : : : :!}()}



{

Return the starting address of the function in the mutatee's address space.}



{

BPatch_flowGraph *getCFG\index{getCFG: : : : : : :!}()



{

Return the control flow graph for the function, or NULL if this information is not available. The BPatch_flowGraph is

described in section \ref{bkm:Ref160277935}.}





\bigskip



{

bool findOverlapping\index{getCFG: : : : : : :!}(std::vector<BPatch_function *> &funcs)



{

Determine which functions overlap with the current function (see Section \ref{bkm:Ref270681825}). Return true if other

functions overlap the current function; the overlapping functions are added to the funcs vector. Return false if no

other functions overlap the current function.}



{

bool addMods\index{getCFG: : : : : : :!}(std::set<StackMod *> mods) \newline

 implemented on x86 and x86-64}



{

Apply stack modifications in mods to the current function; the StackMod class is described in section

\ref{bkm:Ref419102257}. Perform error checking, handle stack alignment requirements, and generate any modifications

required for cleanup at function exit. addMods atomically adds all modifications in mods; if any mod is found to be

unsafe, none of the modifications in mods will be applied.}



{

addMods can only be used in binary rewriting mode.}



{

Returns false if the stack modifications are unsafe or if Dyninst is unable to perform the analysis required to

guarantee safety.}



